\section*{第 \ref{ch:g9:ions} 章——离子与离子溶液}

\begin{solution}{exo:g9:ions:1}
\omterm{def:g9:ions:ion}{离子}是指失去或得到\omterm{def:g9:inside-the-atom:nucleus}{电子}并带有电荷的
\omterm{def:g7:atoms-and-molecules:atom}{原子}或\omterm{def:g7:atoms-and-molecules:atom}{原子}团。\omterm{def:g9:ions:ion}{阳离子}带正电（失去\omterm{def:g9:inside-the-atom:nucleus}{电子}），
\omterm{def:g9:ions:ion}{阴离子}带负电（得到\omterm{def:g9:inside-the-atom:nucleus}{电子}）。
\end{solution}

\begin{solution}{exo:g9:ions:2}
\ce{Mg^{2+}}：12 个\omterm{def:g9:inside-the-atom:proton}{质子}，10 个\omterm{def:g9:inside-the-atom:nucleus}{电子}。
\end{solution}

\begin{solution}{exo:g9:ions:3}
\ce{KCl}、\ce{MgCl2}、\ce{CaO}。
\end{solution}

\begin{solution}{exo:g9:ions:4}
食盐水中含有\omterm{def:g9:ions:ion}{离子} \ce{Na+(aq)} 和 \ce{Cl-(aq)}，它们能移动并
传导电荷。蔗糖\omterm{def:g7:atoms-and-molecules:molecule}{分子}不带电荷。
\end{solution}

\begin{solution}{exo:g9:ions:5}
\ce{CuSO4}、\ce{AlCl3}、\ce{Na2CO3}。
\end{solution}

\begin{solution}{exo:g9:ions:6}
铁（III），\ce{Fe^{3+}}：\ce{Fe^{3+}(aq) + 3OH-(aq) -> Fe(OH)3(s)}。
\end{solution}

\begin{solution}{exo:g9:ions:7}
硫酸根：多 2 个\omterm{def:g9:inside-the-atom:nucleus}{电子}。硝酸根：多 1 个。
\end{solution}

\begin{solution}{exo:g9:ions:8}
与硝酸银生成白色\omterm{def:g9:ions:precipitate}{沉淀}的是 \ce{Cl-}。与氢氧化钠
生成蓝色\omterm{def:g9:ions:precipitate}{沉淀}的是 \ce{Cu^{2+}}。氯化铜（II），\ce{CuCl2}。
\end{solution}

\begin{solution}{exo:g9:ions:9}
四个：每侧各一个（左、右、上、下）。
\end{solution}

\begin{solution}{exo:g9:ions:10}
\ce{Fe^{2+}(aq) + 2OH-(aq) -> Fe(OH)2(s)}。
\end{solution}

\begin{solution}{exo:g9:ions:11}
氯化锌：钠\omterm{def:g9:ions:ion}{离子}与氢氧化钠不产生\omterm{def:g9:ions:precipitate}{沉淀}，
锌\omterm{def:g9:ions:ion}{离子}则生成白色\omterm{def:g9:ions:precipitate}{沉淀}：
\ce{Zn^{2+}(aq) + 2OH-(aq) -> Zn(OH)2(s)}。
\end{solution}

\begin{solution}{exo:g9:ions:12}
\ce{CaCl2}：2000 个氯\omterm{def:g9:ions:ion}{离子}。电荷总和：$1000 \times (+2) + 2000 \times
(-1) = 0$。
\end{solution}

\begin{solution}{pb:g9:ions:1}
\textbf{1.} 铁（II），\ce{Fe^{2+}}：\ce{Fe(OH)2}。

\textbf{2.} 铁（III），\ce{Fe^{3+}}：\ce{Fe(OH)3}。

\textbf{3.} \ce{Fe^{2+}} 多 1 个\omterm{def:g9:inside-the-atom:nucleus}{电子}：$26 - 2 = 24$，
而 \ce{Fe^{3+}} 为 $26 - 3 = 23$。

\textbf{4.} \ce{Fe^{3+}(aq) + 3OH-(aq) -> Fe(OH)3(s)}。

\textbf{5.} 流出时含有铁（II）\omterm{def:g9:ions:ion}{离子}，与\omterm{def:g4:air-a-mixture-of-gases:air}{空气}接触后
转化为铁（III）\omterm{def:g9:ions:ion}{离子}，进而形成橙色固体。

\textbf{6.} \ce{Fe(OH)3}：$(+3) + 3 \times (-1) = 0$。

\textbf{7.} 流出时则不然：铁\omterm{def:g9:ions:ion}{离子}已\omterm{def:g2:mixing-and-dissolving:dissolve}{溶解}，可穿\omterm{def:g3:separating-mixtures:filter}{过
滤}纸。橙色固体形成后则可以，它是一种
\omterm{def:g9:ions:precipitate}{沉淀}，可被滤纸截留。

\textbf{8.} 否：它是\omterm{def:g6:pure-substances-and-mixtures:mixture}{混合物}（水和\omterm{def:g2:mixing-and-dissolving:dissolve}{溶解}的\omterm{def:g9:ions:ion}{离子}）。刚取出时
且澄清，它是\omterm{def:g6:pure-substances-and-mixtures:homogeneous}{均相混合物}。

\textbf{9.} $56 \times 1.67 \times 10^{-27} = \qty{9.35e-26}{kg}$。

\textbf{10.} $\qty{0.30}{mg} = \qty{0.30e-6}{kg} = \qty{3.0e-7}{kg}$。

\textbf{11.} \omterm{def:g9:ions:ion}{离子}与\omterm{def:g7:atoms-and-molecules:atom}{原子}相比，相差两个或三个\omterm{def:g9:inside-the-atom:nucleus}{电子}，
这些\omterm{def:g9:inside-the-atom:nucleus}{电子}的质量与\omterm{def:g9:inside-the-atom:nucleus}{原子核}相比可忽略不计。

\textbf{12.} $3.0 \times 10^{-7} / 9.35 \times 10^{-26} \approx 3.2
\times 10^{18}$ 个铁\omterm{def:g9:ions:ion}{离子}每升。
\end{solution}
